% Physical page 5 onward: the developed commentary of the concise study,
% organized by the questions the three readings answer differently rather than
% by the order of the Mass. Each unit draws on several appointed elements and
% sets the readings' answers beside one another.

\section{The Propers: Detailed Commentary}
\zlabel{triptych:concise:commentary:start}

\subsection{A will none resists, and an adversary who asks}

The Introit and the Offertory put a question to every reading, and the first
reading makes it its own. If no one resists God's will, how does Satan come to
ask for what God's servant does not deserve, and why is it given? Either the
Introit says too much, and the adversary is a second power whose will God must
reckon with, or the Offertory makes God the author of Job's ruin.

St Gregory the Great answers in the preface to his \work{Moralia} on Job with two
wills at work on one man: ``Him the adversary sought indeed to obtain that he
might prove deficient in godliness, but obtained that he might prove a
proficient therein. The Lord in loving-kindness permitted that to be done, which
the devil in his wickedness required'' (Preface~7). His Latin turns on one verb
with two prefixes, \latin{ut deficeret} and \latin{ut proficeret}: the devil
sought that Job might fall short, and obtained that he might go forward. The
petition is the devil's, the permission God's, and the outcome follows the
permission. On the first permission, ``only put not forth thy hand upon his
person'' (Job 1:12), Gregory bids the reader mark ``how He lets go our enemy,
and keeps him in; how He looses, and yet bridles him'', dealing out trials by
seasons, ``that what by coming all together would destroy, may when divided be
borne up against'' (II.19), as St Paul says that God ``will not suffer you to be
tempted above that ye are able'' (1~Cor 10:13). He meets the Introit's claim from
another verse of Job, ``Who hath resisted Him, and had peace?'' (9:4).
Resistance happens, but never succeeds and never has peace: ``Not that our
frailty does resist His unchangeable decree, but what it has not the power to
accomplish, it yet attempts'' (IX.5). Rupert of Deutz, citing the Introit's
words in his treatise on God's omnipotence (he gives them as St Paul's), completes them from that verse and
Gregory's exposition of it: no one resists God's will and has peace (\work{De
omnipotentia Dei} 9).

St Thomas Aquinas states in one distinction what Gregory tells as a story.
Satan harms \latin{non quantum vult sed quantum permittitur}, not as much as he
wills but as much as he is permitted; the will to harm is his own, the power only
from God (\work{Expositio super Iob} 1). In the \work{Summa} he takes the
Introit's own verse from Esther as his authority that nothing can occur outside
the order of the divine government, since God is ``the first universal cause,
not of one genus only, but of all being in general'' (I q.~103 a.~7). Read with
the two Doctors, the Offertory's \latin{petiit} is the adversary's will, his own
and evil; its \latin{data est \dots\ a Domino} is his power, God's and measured;
and the Introit's will is the order to which even the petition belongs.
Gregory's lost peace, Thomas's inescapable order and Bl.\ Ildefonso Schuster's
malice that Providence turns to its own ends, ``Like the devil, who, despitefully
winnowing the grain, renders a helpful service to the great Householder''
(\work{The Sacramentary} III, pp.~178--179), agree on the one point the first
reading needs: the adversary is a creature whose hatred is real and whose reach
is lent. The Collect's \latin{te protegente} then asks for the protection
Gregory's bridle supplies, and Schuster hears in it a prayer against the
obstacles the devil raises to turn men from good works (p.~179).

The other two readings hear the same two chants through other words of them.
The second remembers that Mardochai prays for a people under sentence, and in
lines the Mass does not sing asks, ``have mercy on thy people'' (Esth 13:15); it
meets in the Offertory's Job the man whom St John Chrysostom, preaching on this
Mass's parable, sets before the hearer as proof that even the devil's malice can
become a cause of crowns (\work{Homilies on Matthew} 61.5). The third hears
Mardochai remind God of ``thy portion, which thou hast redeemed for thyself out
of Egypt'' (13:16), and reads Job with Gregory, who takes the name to mean
``grieving'' and to set forth ``either our Mediator's Passion, or the travails
of Holy Church, which is harassed by the manifold toils of this present life''
(Preface~16); Schuster says that the history of Job ``is repeated again and
again throughout the history of the Church'' (p.~179). Two liturgical
commentators whom the first reading follows hear in the Offertory a soldier: for Rupert what happened
to Job in history happens in the moral life to every soldier of the virtues,
whom the tempter strikes when a hidden dispensation of God permits him (\work{De
divinis officiis} XII.21), and for William Durandus the devil is conquered above
all by patience (\work{Rationale} VI.138).

The objection that this makes God the author of Job's ruin is answered by
keeping the two wills apart. The will to harm is the devil's and evil; the
permission is God's and good; and God grants what the devil asks only so far as
it serves the proving of his servant. \latin{Ut deficeret} and \latin{ut
proficeret} carry the whole answer: the same trial is the devil's attempt at
Job's fall and God's means of his advance.

\subsection{Who is the enemy, and what is the armour?}

The Epistle names the enemy, and Chrysostom presses the answer hard. Paul says
``wiles'' because the devil ``never proposes to us sins in their proper
colors''. The rulers of this darkness are incorporeal, and the Christian
overcomes them by becoming what they are not: ``How then, you may say, are we to
wrestle with the darkness? By becoming light.'' Then the decisive turn. Paul
wrote to disciples who would have many persecutors, and told them in effect,
``imagine not \dots\ that it is they who war with you. They that really war with
you, are the spirits that work in them'' (\work{Homilies on Ephesians} 22). From
this Chrysostom draws a rule stated with deliberate violence: let the man who has
been wronged ``collect together all that wrath, all that fierceness, and pour it
out upon the head of the devil. Here wrath is a good thing, here anger is
profitable, here revenge is praiseworthy.'' The wrath is not abolished but
redirected, from the brother to the adversary who moved him. St Thomas, answering
the objection that men plainly do persecute us, allows that they do, but as moved
by a higher agent, \latin{a superiore movente, sc.\ a diabolo} (\work{Super
Epistolam ad Ephesios} 6, lect.~3). He agrees that the spirits are the real
enemy, and by calling men instruments he keeps the persecutor answerable for
what he does. Chrysostom too tells the wronged man where to turn his wrath, not that the
wrong was none, and his praiseworthy revenge is aimed at the devil alone.

Here the first two readings meet at a single thought of Chrysostom's, which he
states in the homily on the parable and joins to the Offertory's Job: ``But if
the devil hath become a cause of crowns, why art thou afraid of a man as an
enemy?'' (\work{Homilies on Matthew} 61.5). The first reading takes its first
half: the enemy is the adversary, not the brother. The second takes the other
half: the brother is the fellow servant whose small debt the forgiven man must
forgive. The two are complementary and distinct. The first answers a question
about God's government and would stand without the parable; the second answers
a question about the forgiven man's heart and would stand without the
Offertory's permission. The first reading even finds the adversary at the
Gospel's end, where St Thomas glosses the torturers as the demons to whom the
unforgiving servant is subjected, \latin{subiicitur daemonibus} (\work{Super
Matthaeum} 18). On the gospel of peace with which the feet are shod, Chrysostom
adds the paradox by which the first reading sets the Secret's offering within the
war: ``if we are at war
with the devil, we are at peace with God'' (\work{Homilies on Ephesians} 24).

The armour is \latin{armatura Dei}, and the witnesses name it differently. St
Jerome makes it Christ himself: to put on the whole armour of God is to put on
the Lord, \latin{quasi dixerit, Induite Dominum Jesum Christum} (Rom 13:14), so
that girdle, breastplate and shield are each Christ; and he reads the wars of
Israel, David's against Goliath and those of Josue, the Judges and the Kings, as
written for us, a figure of the spiritual war (\work{In Ephesios} III, PL~26,
col.~543). St Thomas names the pieces as the gifts and virtues, the girdle
temperance, the breastplate justice and the shield faith. The continuation of
\work{The Liturgical Year}, commenting on this Mass, finds the same pieces worn by
Christ in the prophets before they are ours: ``well, yes, but they are his,
first'' (XI, p.~460). The first reading holds Jerome's identification and
Thomas's together, as the Christological and the moral sides of one gift: the
soldier is clothed in Christ, and Christ in him becomes virtue. The third reading
takes from the same Epistle not the armour but the girded loins.

The Gradual gives the soldier his fortress. St Robert Bellarmine explains that
the Hebrew for refuge ``signifies a well fortified house, placed on an
eminence'', and that those who dwell in God dwell ``in a city fortified by
faith, hope, and charity'' (\work{Explanation of the Psalms}, Ps~89:1).
Durandus, expounding the Gradual, says that because even the well armed cannot
perfectly resist unless God helps them, the Church flees to God her refuge
(VI.138). The refuge that is God before the mountains is the same will that the
Introit says none resists: what threatens the soldier lies inside it, and so
does he.

\subsection{What the king's mercy requires}

The second reading hears the whole Mass from the Gospel, and everything in the
parable depends on the disproportion of the two debts. A denarius is the day's
wage of another parable's labourers, ``a penny a day'' (Mt 20:2), so the fellow
servant owes about a hundred days' labour, a real debt that a working man could
in time repay. St Augustine supposes the ten thousand talents to be ``at least
ten thousand sins'', and asks of the hundred denarii, ``But what was this to ten
thousand talents?'' (\work{Sermon} 83.3). Chrysostom draws the lesson: ``Seest
thou how great the difference between sins against man and against God?'' It
arises from the difference of the persons and from ``the constant succession of
our sins'', since we shrink from sinning when a man looks at us but not when God
sees us every day (\work{Homilies on Matthew} 61.1). Augustine makes every
hearer both debtors at once: ``every man is at once God's debtor, and hath also
some brother a debtor to himself'' (83.2). Jerome puts the argument as a single
question from the greater to the lesser, \latin{quanto magis servi conservis suis
debent minora dimittere?} (\work{In Matthaeum} III, PL~26, col.~133).

Chrysostom reads the king's first sentence as mercy. The order to sell the
servant with his family came ``of unspeakable tenderness'': ``it is His purpose to
alarm him by this threat, that He might bring him to supplication, not that he
should be sold.'' The servant asked only for time, and the lord ``gave more than he
asked, remission and forgiveness of the entire debt'' (61.3). When the fellow
servant falls down and says the same words, \latin{Patientiam habe in me, et omnia
reddam tibi}, the forgiven man hears his own prayer and refuses it; he ``did not
recognize so much as the harbor by which he escaped shipwreck'' (61.4).

Then the difficulty. The king had said \latin{omne debitum dimisi tibi}; now he
exacts \latin{universum debitum}. Can a pardon return? The witnesses answer in
three formulations that the second reading keeps side by side. Chrysostom
corrects his own phrase that the lord was ``on the point of revoking his gift'':
``or rather, it was not he that revoked it, but the one who had received it''.
The delivery to the torturers is ``for ever; for he will never repay'', and
wickedness has set aside even the law that God's gifts are without repentance
(61.4; Rom 11:29). Augustine puts the sentence in God's mouth: ``For if thou wilt
not forgive, I will call thee back, and put upon thee again all that I had
remitted to thee'' (83.7). For Chrysostom the receiver revokes the gift; for
Augustine God calls it back. The two are complementary within their own grammar,
the one looking at the cause and the other at the judge who ratifies it. Jerome
calls the closing words \latin{formidolosa sententia}, fearful if God's sentence
bends and changes according to our own mind, and says that it overturns every
pretence of a feigned peace, \latin{omnem simulationem fictae pacis evertit}
(col.~133). St Thomas answers more narrowly, and here the difference is real. The
verse seems to suggest that forgiven sins return, as Origen held; Thomas denies
it, since the remission of sins has its efficacy from the sacraments, and says
that they are said to return \latin{per ingratitudinem}, through ingratitude
(\work{Super Matthaeum} 18). He does not deny the servant's fate; he denies that
the sins once remitted come back as the same sins. Thomas tells the hearer that a
grudge does not undo his baptism; Augustine and Chrysostom tell him that refusing
mercy after receiving it makes him again a debtor of the whole. Each keeps the
other from being misheard.

The parable's last words, \latin{de cordibus vestris}, carry the condition.
Chrysostom finds two things required, ``both to condemn ourselves for our sins,
and to forgive others'', the former for the sake of the latter, ``and not merely
to forgive with the lips, but from the heart'' (61.5). Augustine: ``Say not with
the tongue, `I forgive,' and put off to forgive in the heart'', and, reversing
the ordinary preference, ``Better is it that thou shouldest be violent in words,
and forgive in the heart, than in words be soft, and in the heart relentless''
(83.7). He also shows what the words do not forbid. A boy who sins, asks forgiveness and sins again may still be
corrected, ``with words, and if need be with scourges'', so long as the sin is
forgiven and its remembrance cast from the heart; ``therefore did the Lord add,
`from your hearts,' that though through affection discipline be exercised,
gentleness might not depart out of the heart'' (83.8). Forgiveness from the heart
forbids rancour, not correction. And Augustine names where the promise is made
daily, ``Forgive us our debts, as we forgive our debtors'': ``this the covenant
and agreement thou dost mention when thou prayest'' (83.4). Schuster calls the
parable ``the most authoritative commentary on those words of the Lord's
Prayer'', and adds that ``The measure of divine compassion is fixed by
ourselves'' (p.~180).

Who is the debtor? Some, Jerome reports, made him the devil, and he rejects the
reading outright as neither the Church's nor one for prudent men (col.~133).
Chrysostom and Augustine make him every hearer, and Augustine says: ``every man begins from Baptism; he goes out
free, the `ten thousand talents' are forgiven him; and when he goes out, he will
soon find some fellow-servant his debtor'' (83.7). St Thomas makes the king God
and the servants every man, while St Rabanus Maurus, in an allegory of his own,
reads the debtor as the people of the Jews bound by the ten commandments and the
fellow servant as the Gentiles, before he too turns to every hearer. The
parable's \latin{unusquisque fratri suo} settles the moral identification: the
debtor is each of us, and the fellow servant each brother.

The second reading then hears the other elements surround the parable. The
Introit's undefiled are forgiven debtors: Augustine will not let them forget
``If we say that we have no sin, we deceive ourselves'' (\work{Enarrationes}
118, sermon~1.3), and Cassiodorus says that the saints become undefiled
\latin{quando indulgentiae munera consequuntur}. The Collect asks for the whole
household the \latin{continua pietas} the king showed his servant, and its good
works include the two that Augustine sets beside the Gospel, ``Forgive, and ye
shall be forgiven: give, and it shall be given unto you'' (83.2). In the parable
the lord was appeased by a prayer, and Chrysostom reminds the servant of ``God,
whom by mere prayer thou hast reconciled'' (61.4); the Secret offers the
sacrifices by which God willed to be appeased, and Schuster reads its
\latin{potenti pietate} as the mercy that ``has triumphed over divine justice
itself'' (p.~181). The Postcommunion asks that what the mouth has received the
pure mind may follow; Chrysostom's ``not merely \dots\ with the lips, but from
the heart'' asks the same agreement of mouth and heart of the one who forgives,
and so the Mass ends where the parable ends.

\subsection{When wilt thou judge? The Communion's persecutors}

The plea \latin{quando facies de persequentibus me iudicium?} is where the three
readings differ most, and the difference runs between the witnesses themselves.
Each reading must meet the objection that it asks for vengeance after a Gospel
of forgiveness, and the second answers it from Augustine in two places. In a sermon on anger: ``A man
lusteth to be avenged, and Christ is not yet avenged, the holy martyrs are not yet
avenged. Still doth the patience of God wait, that the enemies of Christ, the
enemies of the martyrs, may be converted. \dots\ Forgive therefore; from the heart
forgive'' (\work{Sermon} 58.8). And of the martyrs who cry to be avenged he asks
``whether they pleaded against the men themselves or against the dominion of
sin'', and answers that ``it is a genuine avenging of the martyrs, and one full of
righteousness and mercy, that the dominion of sin should be overthrown''. It is
overthrown partly by men's amendment and partly by the condemnation of the
incurable; St Paul, bringing his body into subjection, avenged Stephen ``in his
own person'' (\work{De sermone Domini in monte} I.22.77). For the second reading
the plea and the command ask the same thing, the end of sin's rule, whose first
way of ending is the persecutor's conversion.

Who, then, are the persecutors, and when is the judgment? St Hilary of Poitiers
names them in the Epistle's terms: \latin{a diabolo scilicet, ministrisque ejus.
Sunt enim spiritales nequitiae in caelestibus}. What the singer asks is that,
made wholly holy, he may crush the devil within himself, \latin{conterat intra se
diabolum}; it is the power the Lord gave his apostles ``to tread upon serpents
and scorpions, and upon all the power of the enemy'' (Lk 10:19; \work{Tractatus
in Psalmum} 118, \latin{Caph}~5, PL~9, cols.~573--574). Hilary knows the
judgment after the end of the world, and that the devil is to be judged at the
resurrection, but says that what the prophet now asks is \latin{non illud utique
quod in resurrectionis saeculo meminit futurum}. His ``now'' rests on the
question the Communion does not sing, ``How many are the days of thy
servant?'': the prophet knows himself a man of few days, and asks for what he may
have in them. The first reading follows him: the judgment asked for is not a
curse on men but the power to tread the adversary down.

Augustine hears another voice. ``In the Apocalypse, these are the words of the
Martyrs, and long-suffering is enjoined them until the number of their brethren
be fulfilled'' (\work{Enarrationes} 118, \latin{Caph}~84): the souls under the
altar who cry ``How long, O Lord \dots\ dost thou not judge and revenge our blood''
and are told to rest ``till their fellow servants and their brethren \dots\ should
be filled up'' (Apoc 6:10--11). Reading the verse's two questions in order, he
hears first that ``The body of Christ then is asking concerning its days, what
they are to be in this world'', and then the judgment, ``showing indeed that the Church would exist on earth until the
judgment, when vengeance shall fall upon Her persecutors''. That ``until'' comes
from the order of the two questions, the first of which the Communion leaves
unsung; on the half that is sung he gives the martyrs' plea answered with
long-suffering. On the cry for help he adds its object, ``that I may struggle for
the truth even unto death'', against men who persecuted ``nothing save the
truth'' (\latin{Caph}~86). Cassiodorus hears a people on pilgrimage asking for the
end of the world, so that the vengeance prepared for the persecutors, \latin{id est
diabolo cum ministris}, may be completed, \latin{non quia vitae suae tempus volebat
agnoscere: sed celeritatem futuri iudicii postulabat} (\work{Expositio Psalmorum}
118:84). Bellarmine reads the sung half the same way: ``When will you judge and
condemn those that persecute and tempt me \dots\ He makes no curious inquiries as to
the number of the days of his life'' (\work{Explanation of the Psalms}, Ps~118:84).
The exposition of the psalms printed under St Athanasius's name holds both
answers: on the unsung question it asks that Satan be crushed at once under the
singer's feet (Rom 16:20), and on the sung cry it says that the man persecuted
unjustly, knowing he will receive his recompense at the divine tribunal,
\latin{gestit illud tempus videre}, longs to see that time (PG~27, col.~495).

The witnesses differ, then, in the thing asked, a present power or the last
judgment; in the persecutors named, the devil with his ministers for Hilary and
Cassiodorus, men for Augustine and Bellarmine; and in the time, which each draws
from the unsung question. The answers are compatible without being the same.
Hilary affirms the judgment that Augustine and Bellarmine await, and the men who
persecute are, in St Thomas's terms, instruments moved by the devil. The first
reading follows Hilary's identification and the power he asks for; the second,
Augustine's dominion of sin; the third, the end-directed answer of Augustine,
Cassiodorus and Bellarmine. The liturgical commentators hear first the hope in
which the question is asked: for Durandus patience avails nothing without hope
(VI.138), and for the continuation of \work{The Liturgical Year} ``An unflagging
hope ever accompanies the admirable patience of holy Church'' (p.~470).

\subsection{A people on the march}

The third reading gives the waiting a road, and finds it in the Alleluia and the
Epistle together. Augustine teaches that ``no Christian consider himself alien to
the name of Israel''; that Egypt ``is often used for an emblem of this world;
from which we must spiritually withdraw''; and that ``as they did not depart
thence, until freed by Divine help; so no man is turned away in heart from this
world, unless aided by the gift of the Divine mercy''. The going out ``is fulfilled
in every faithful one in the daily travailings of the Church'' (\work{Enarrationes}
113.3--4). Cassiodorus says that we truly become Israel \latin{cum pompae mundi
huius a nostra coeperint mente discedere}, when the pomps of this world begin to
leave our minds; he also reads the going out as release from the host of the
Egyptians, \latin{id est a daemonum turba}, which is how the first reading hears
the verse (\work{Expositio Psalmorum} 113.1). The second, with Augustine, hears in
both releases the work of mercy.

Chrysostom finds the same Exodus in the Epistle's \latin{State ergo succincti
lumbos vestros}. Israel ate the Passover girded, shod and with staves in hand
because ``we were ready for our exodus''. ``This then is the historical type. But
the reality is this; we too eat a Passover, even Christ \dots\ We too ought to eat
it, both sandalled and girded. And why? That we too may be ready for our Exodus,
for our departure hence.'' The desert lies ahead, and the devil in it; but ``We
have manna; yea rather not manna, but far more than manna'', and ``the leader of
this our Exodus is not Moses, but Jesus''. No one builds a splendid house in an
encampment, for ``The present life is nothing else than a march and an
encampment'' (\work{Homilies on Ephesians} 23). The Epistle's \latin{dies malus}
then holds two meanings on one road. Chrysostom: ``By `evil day' he means the
present life'' (Hom.~22). The continuation of \work{The Liturgical Year}: evil
days are many, ``But, for every man, and for the world at large, there is one
evil day, evil beyond all the others: it is the last day, the day of judgment''
(XI, p.~459). The whole march is an evil day, and its end is the day that ends
the others.

The Gradual names the God toward whom the march goes. Augustine marks the
present tense of \latin{tu es, Deus}: ``Behold then the eternity that is our
refuge, that we may fly thither from the mutability of time'' (\work{Enarrationes}
89.3). Where the first reading heard a fortress, the third hears an eternity, and
Fromage hears in the Gradual and the Alleluia together the Church ``preparing the
last exodus of the Church, which is leaving this world turned infidel, and is
hastening to the true land of promise'' (p.~463).

The Gospel's king gives the march its end, and here the second and third readings
share a witness while placing the weight differently. St Thomas says that the
reckoning begins when God first sends some tribulation and ends at the day of
judgment, \latin{Finis huius rationis erit in die iudicii; principium quando primo
inducit aliquam tribulationem} (\work{Super Matthaeum} 18). Honorius
Augustodunensis and Durandus, expounding the parable in their books, read the
reckoning of the judgment (\work{Gemma animae} IV.93; \work{Rationale} VI.139),
and Fromage sets beside this Gospel the \latin{Dies irae}'s petition, ``grant us
the gift of forgiveness, before the Day of reckoning cometh'', saying, ``We are all
of us, in fact, that negligent servant, that insolvent debtor'' (pp.~464--466).
For the second reading the reckoning is the condition's sanction: St Alphonsus
Liguori calls \latin{tradidit eum tortoribus} ``a description of the sentence of
the eternal death'' (\work{Sermons for All the Sundays}, Sermon~L), agreeing with
Chrysostom's ``for ever''. For the third it is the destination: the tribulations
of the march are the reckoning's beginning, and the judgment the Communion asks
for is its end.

The rest of the Mass travels the same road. Mardochai asks, in lines the Mass does
not sing, ``that we may live and praise thy name, O Lord'' (Esth 13:17), and
Fromage reads that petition as the reason for the Collect's prayer for peace: the
Church asks freedom from adversity so that she may pay God ``the united homage of
religion and good works'' (p.~456). The Introit's verse blesses those who walk,
\latin{qui ambulant}. The Secret asks that salvation be restored, \latin{salutem
\dots\ restitui}, so that the salvation the Communion hopes in is given on the way
as well as at its end. And the Postcommunion names the food: \latin{immortalitatis
alimonia}. Chrysostom's Passover that is Christ, ``far more than manna'', is said
of the Epistle's girded loins; the Postcommunion names the same food in its own
words, the nourishment of those who will not die, eaten by a people still on the
march. Bellarmine's gloss on the Communion's question, ``his desire of life
everlasting'', names what that food answers.

\subsection{One Mass, three answers}

Each reading takes up a demanding feature of the formulary that the others leave
in the background. The first takes the Introit's absolute sentence beside the
Offertory's \latin{petiit} and \latin{data est \dots\ a Domino}, and answers it
with Gregory's two wills and Thomas's distinction of will and power. The second
takes the Gospel's reversed pardon and its \latin{universum debitum}, and answers
without softening the sentence or letting it unsettle the efficacy of God's
forgiveness. The third takes the Communion's \latin{quando} and answers by the
time it opens, the delay of the judgment in which God's patience waits for the
persecutors' conversion. The second and third meet at Augustine's sermon on
anger, where the martyrs not yet avenged and the command to forgive from the heart
stand in one sentence: the second uses it to show that the plea does not undo the
Gospel, the third to show that the waiting is the time of God's patience. Both
hold that the Church that asks for judgment forgives in the meantime.

Their four senses differ where their questions differ. In the allegorical sense
the first finds Christ as the armour and the Church delivered from the devil as
the people delivered from Aman; the second, God as the king and baptism as the
remission of the ten thousand talents; the third, the Church's exodus out of the
world with Christ for its Passover. The moral senses ask three distinct things:
stand, and turn the wrath from the brother to the devil; forgive every debt from
the heart, and correct without rancour; live as on a march, and forgive while
God's patience waits. In the anagogical sense the first looks to the adversary
judged at the resurrection and the soldier crowned; the second, to the reckoning
completed, where God's measure will be the measure we used and the Church asks
beforehand for the gift of forgiveness; the third, to the land of promise toward
which the last exodus hastens, fed on the food of immortality. Read together, the
three show one Mass in which the God who bounds the adversary is the king who
forgives the debt, and the people that stands in the evil day is the people that
goes out of Egypt toward his judgment.
